\subsection{Rational normalization in characteristic two}

For a very ample curve in characteristic two, a positive-genus ruling is incompatible with the Euler characteristic and intersection form. The following proposition reduces the smallest equality to a rational normalized cover; the canonical argument and Theorem~\ref{thm:strict-picard-all} below exclude that case.

\begin{proposition}\label{prop:characteristic-two-rational-normalization}
Let $X$ be a smooth projective integral surface over an algebraically
closed field $k$ of characteristic two. Let $Y\subset X$ be a smooth
integral very ample divisor, and let $M\ne\OO_X$ have order two with
$M|_Y\simeq\OO_Y$. Suppose that
\[
 2g(Y)-2=2Y^2.
\]
Let $\pi:N\to X$ be the normalization of the $\mu_2$-torsor defined by a
trivialization of $M^2$. Then $N$ is birationally rational. In particular,
$X$ admits a dominant purely inseparable rational map of degree two
from a rational surface. The surface $X$ itself is not birationally
ruled.
\end{proposition}

\begin{proof}
Corollary~\ref{cor:very-ample-equality-almost-ruled} and
Proposition~\ref{prop:distinguished-ruled-cover} give a presentation
$k(N)=k(C)(t)$, with $C$ a smooth projective integral curve, and show
that $X$ is not birationally ruled. We prove that $g(C)=0$.

Suppose instead that $g=g(C)\geq1$, and put $d=Y^2>0$.
Adjunction gives
\[
 K_X\cdot Y=d.
\]
Proposition~\ref{prop:positive-genus-picard-fibration} supplies the
morphisms
\[
 f:X\longrightarrow B=C^{(2)},\qquad h:N\longrightarrow C,
 \qquad f\pi=F_{C/k}h,
\]
and gives
\[
 \chi(\OO_N)=1-g.
\]
For a general integral fibre $F$ of $f$, put $b=Y\cdot F$.
The same proposition gives
\[
 F^2=0,\qquad K_X\cdot F=b-2,\qquad
 2p_a(F)=b,\qquad b\geq4.
\]
In particular $b$ is even. The fibration and fibre identities here
require only ampleness once the ruled cover is given; very ampleness
is used for the last bound.

We first compute the Euler characteristic of the finite flat
rank-two algebra $\pi_*\OO_N$. Its unit is nonzero on every fibre,
so the quotient by $\OO_X$ is an invertible sheaf $Q$.
The inclusion of the Kummer algebra induces $M^{-1}\to Q$.
At codimension-one points outside $Y$ this is an isomorphism.
At the generic point of $Y$, the conductor basis is $1,w$, with
$w=(z-v)/t$, where $t$ is a local equation of $Y$. Modulo the unit
summand, $[z]=t[w]$. The divisor of the map of line bundles is
therefore exactly $Y$, giving
\[
 0\longrightarrow\OO_X\longrightarrow\pi_*\OO_N
 \longrightarrow M^{-1}(Y)\longrightarrow0;
\]
this is also the first assertion of
Proposition~\ref{prop:k3-picard-paracanonical}.
Since $M$ is numerically trivial, Riemann--Roch
\cite[Tag~02UO]{Stacks} and
$K_X\cdot Y=Y^2$ yield
\[
 \chi(M^{-1}(Y))
 =\chi(\OO_X)+\tfrac12Y\cdot(Y-K_X)=\chi(\OO_X).
\]
Finite pushforward and additivity of Euler characteristic thus give
\begin{equation}\label{eq:characteristic-two-cover-euler}
 2\chi(\OO_X)=\chi(\OO_N)=1-g.
\end{equation}
No splitting of the rank-two algebra is required.

For the classical adjunction argument used below, one may first extend
$k$ to an uncountable algebraically closed field and perform the
whole argument there. All the geometric hypotheses and the genus
of $C$ persist. The line bundle $M$ remains nontrivial: a nonzero
section of a torsion line bundle would have an effective divisor
with zero ample intersection and hence would trivialize it;
thus $H^0(X,M)=0$, and this vanishing persists under field extension.
The divided differential and
Proposition~\ref{prop:equality-rational-curves} again exclude
$(-1)$-curves after extension, and
Proposition~\ref{prop:distinguished-ruled-cover} again excludes
birational ruledness. The ruled blowup model and its degree-two
map also base change. A contradiction over the enlarged field
therefore suffices; no descent of ruledness is needed.

Proposition~\ref{prop:equality-rational-curves} excludes every
$(-1)$-curve on $X$. The classical adjunction-termination theorem
therefore makes $K_X$ nef: a negative canonical intersection on a
smooth projective surface without $(-1)$-curves would make that
surface birationally ruled
\cite[\S1, Steps~I--VII, pp.~326--330]{MumfordEnriquesI}.

Here $K_X^2\geq0$ can be checked directly from nefness and
$K_X\cdot Y=d>0$. If $K_X^2=-q<0$, then for rational $t>d/q$
sufficiently close to $d/q$, the class $L=Y+tK_X$ satisfies
\[
 L^2>0,\qquad L\cdot Y>0,\qquad L\cdot K_X<0.
\]
After clearing denominators, take $L$ integral.
For large $n$, Serre duality gives
\[
 H^2(X,\OO_X(nL))=0,
\]
because $(K_X-nL)\cdot Y<0$. Riemann--Roch then gives a nonzero
effective divisor in $|nL|$, contradicting the negative
intersection $K_X\cdot nL$. Thus $K_X^2\geq0$.

Fix a prime $\ell\ne2$, and write $b_i(X)$ for the
$\ell$-adic Betti numbers. The ruled model furnished by
Proposition~\ref{prop:distinguished-ruled-cover} is a smooth surface
$S$ obtained by point blowups of $C\times\PP^1$, with a purely
inseparable morphism $S\to X$ of degree two. Ekedahl's
Betti-number lemma gives
\[
 b_1(X)=b_1(S)=2g
 \quad\text{\cite[Lemma~0.1.5, p.~103]{EkedahlCanonical}}.
\]
The lemma does not require $S\to X$ to be finite. Its proof uses
injectivity of pullback in cohomology and a reverse map after
a Frobenius twist and point blowups.

The classes of $Y$ and $F$ in
$H^2_{\mathrm{\acute et}}(X,\mathbf Q_\ell(1))$
have intersection matrix
\[
 \begin{pmatrix}d&b\\ b&0\end{pmatrix}.
\]
Its determinant is $-b^2\ne0$, so the two classes are independent
and $b_2(X)\geq2$. Poincar\'e duality and algebraic Gauss--Bonnet
\cite[Tag~0FH0]{Stacks} give
\[
 c_2(X)=2-2b_1(X)+b_2(X)
       =2-4g+b_2(X)\geq4(1-g).
\]
On the other hand, Noether's formula and
\eqref{eq:characteristic-two-cover-euler} give
\[
 c_2(X)=12\chi(\OO_X)-K_X^2\leq6(1-g).
\]
For $g>1$ these two inequalities contradict each other.

It remains to consider $g=1$. The same formulas force
\[
 \chi(\OO_X)=0,\qquad c_2(X)=0,\qquad
 K_X^2=0,\qquad b_2(X)=2.
\]
The independent classes of $Y$ and $F$ now form a basis of
$H^2_{\mathrm{\acute et}}(X,\mathbf Q_\ell(1))$.
Write $[K_X]=\alpha[Y]+\beta[F]$. Pairing with $F$ and $Y$
gives
\[
 \alpha b=b-2,\qquad \alpha d+\beta b=d,
 \qquad
 \alpha=1-\frac2b,\qquad \beta=\frac{2d}{b^2}.
\]
Compatibility of the cup product with intersection then yields
\begin{equation}\label{eq:characteristic-two-canonical-square}
 K_X^2=\alpha^2d+2\alpha\beta b
      =d\left(1-\frac4{b^2}\right)>0,
\end{equation}
because $b\geq4$. This is the final contradiction.

Hence $g(C)=0$. Over the algebraically closed field, $C\simeq\PP^1$,
so $k(N)$ is a rational function field in two variables.
The extension $k(N)/k(X)$ is purely inseparable of degree two;
a rational model of $N$ supplies the asserted rational domination
of $X$.
\end{proof}

\begin{remark}\label{rem:characteristic-two-ample-residual}
Suppose that $Y$ is merely ample, with the same characteristic-two
equality and torsion hypotheses, and additionally suppose that $X$
has a purely inseparable ruled rational cover of degree two.
If the resulting ruling curve $C$ has positive genus, the proof
above still gives
\[
 g(C)=1,\qquad
 \chi(\OO_X)=K_X^2=c_2(X)=0,\qquad b_1(X)=b_2(X)=2.
\]
The general-fibre identities give a positive even integer $b=Y\cdot F$.
Equation~\eqref{eq:characteristic-two-canonical-square}, before its
strict inequality, then forces $b=2$ and $p_a(F)=1$.
Moreover, $Y\to B=C^{(2)}$ is separable of degree two. Its degree is
$Y\cdot F=2$; if it were inseparable, it would be a relative
Frobenius followed by an isomorphism and would preserve genus.
That would give $g(Y)=1$, contrary to $g(Y)=Y^2+1\geq2$.
Thus the purely inseparable multisection obstruction in
Proposition~\ref{prop:quasielliptic-multisection} does not already
exclude this configuration.

The genus-zero branch for merely ample $Y$ satisfies the necessary
numerical conditions in Proposition~\ref{prop:rational-normalization-canonical}.
The present proposition concerns the function field of the normalized
cover and does not extend a chosen rational ruling map over
every point of that cover. It supplies no existence example for
the Picard kernel, and does not settle the higher-conductor
or strict-inequality cases.
\end{remark}
